%%%%%%%%%%%%%%%%%%%%%%%%%%%%%%%%%%%%%%%%%
%%%%%%%%%%%%%%%%%%%%%%%%%%%%%%%%%%%%%%%%%
\section{Optimization Methodology}\label{remdo:sec:optim-methods}
\ifdefined\DISSERTATION
The methodology is iterative and consists of four main components: 
(1) decide the simulation structure, 
(2) implement and verify appropriate models for each simulation module, 
(3) formally define the optimization problem formulation, and 
(4) outline a process for solving the optimization problem and analyzing the results.
Figure~\ref{remdo:fig:overview-methods} depicts this methodology visually.
\else
The methodology is iterative and consists of four main components: 
(1) decide the simulation structure, 
(2) implement and verify appropriate models for each simulation module, 
(3) formally define the optimization problem formulation, and 
(4) outline a process for solving the optimization problem and analyzing the results.
Figure~\ref{remdo:fig:overview-methods} depicts this methodology visually.
\fi
\begin{figure}
    \centering
    \includegraphics[width=1\linewidth]{figs/from-matlab/methodology_flowchart.pdf}
    \caption{Methodology overview}
    \label{remdo:fig:overview-methods}
\end{figure}
\subsection{Simulation Structure}
The optimization uses a simulation model known as MDOcean, a purpose-developed open-source MATLAB package for fast and accurate semi-analytical WEC modeling.
\input{sections/model-overview}

\subsection{Modeling Approach}\label{remdo:sec:model}
This section provides an overview of the modeling approach, while readers are referred to reference \cite{mccabe_development_2026} for the full assumptions and validation of the model.
%\hl{Todo: summarize core equations here - one per module}
\ifdefined\DISSERTATION
\else
\fi

\begin{figure}
    \centering
    \includegraphics[width=\linewidth]{figs/from-matlab/xdsm.pdf}
    \caption{Simplified extended design structure matrix (XDSM) diagram of the MDOcean modules, adapted from \cite{mccabe_development_2026}.
Inputs enter at the top, data passes forward between modules left to right along the diagonal, and outputs exit at the left; the orange block denotes the internally iterated dynamics module.}
    \label{remdo:fig:n2}
\end{figure}

\ifdefined\DISSERTATION
\else
\fi
